\documentclass[10pt,a4paper]{ctexart}

\usepackage[a4paper,margin=1.45cm,headheight=14pt,footskip=17pt]{geometry}
\usepackage{fontspec}
\usepackage{xeCJK}
\usepackage{xcolor}
\usepackage{booktabs}
\usepackage{tabularx}
\usepackage{array}
\usepackage{enumitem}
\usepackage{hyperref}
\usepackage{fancyhdr}
\usepackage[most]{tcolorbox}
\usepackage{amsmath}
\usepackage{microtype}

\setCJKmainfont{Noto Sans CJK SC}
\setCJKsansfont{Noto Sans CJK SC}
\setmainfont{Noto Sans CJK SC}
\setsansfont{Noto Sans CJK SC}
\setmonofont{DejaVu Sans Mono}

\definecolor{navy}{HTML}{17324D}
\definecolor{blue}{HTML}{2563EB}
\definecolor{bluebg}{HTML}{EFF6FF}
\definecolor{green}{HTML}{0F766E}
\definecolor{greenbg}{HTML}{ECFDF5}
\definecolor{amber}{HTML}{B45309}
\definecolor{amberbg}{HTML}{FFFBEB}
\definecolor{red}{HTML}{B91C1C}
\definecolor{redbg}{HTML}{FEF2F2}
\definecolor{slate}{HTML}{52657A}
\definecolor{line}{HTML}{CBD5E1}

\hypersetup{
  colorlinks=true,
  linkcolor=navy,
  urlcolor=blue,
  pdftitle={Agent Skill Acquisition：文献与研究进展汇报},
  pdfsubject={导师私聊汇报与组会选篇},
  pdfauthor={Agent Skill Acquisition Research}
}

\pagestyle{fancy}
\fancyhf{}
\fancyhead[L]{\scriptsize Agent Skill Acquisition}
\fancyhead[R]{\scriptsize 文献与研究进展 · 2026-07-24}
\fancyfoot[C]{\scriptsize\thepage\,/\,4}
\renewcommand{\headrulewidth}{0.35pt}

\setlength{\parindent}{0pt}
\setlength{\parskip}{0.22em}
\setlist[itemize]{leftmargin=1.35em,itemsep=0.09em,topsep=0.15em,parsep=0pt}
\setlist[enumerate]{leftmargin=1.55em,itemsep=0.09em,topsep=0.15em,parsep=0pt}
\renewcommand{\arraystretch}{1.12}

\newcolumntype{Y}{>{\raggedright\arraybackslash}X}
\newcolumntype{L}[1]{>{\raggedright\arraybackslash}p{#1}}
\newcolumntype{C}[1]{>{\centering\arraybackslash}p{#1}}
\newcommand{\arxiv}[1]{\href{https://arxiv.org/abs/#1}{#1}}
\newcommand{\pageheading}[2]{%
  {\Large\bfseries\color{navy}#1}\hfill{\small\color{slate}#2}\par
  \vspace{0.12cm}\color{line}\hrule\color{black}\vspace{0.28cm}
}

\newtcolorbox{keybox}{
  enhanced,colback=greenbg,colframe=green,boxrule=0.75pt,arc=1.4mm,
  left=2.4mm,right=2.4mm,top=1.7mm,bottom=1.7mm
}
\newtcolorbox{warnbox}{
  enhanced,colback=amberbg,colframe=amber,boxrule=0.7pt,arc=1.4mm,
  left=2.4mm,right=2.4mm,top=1.6mm,bottom=1.6mm
}
\newtcolorbox{redbox}{
  enhanced,colback=redbg,colframe=red,boxrule=0.7pt,arc=1.4mm,
  left=2.4mm,right=2.4mm,top=1.6mm,bottom=1.6mm
}
\newtcolorbox{bluebox}{
  enhanced,colback=bluebg,colframe=blue,boxrule=0.7pt,arc=1.4mm,
  left=2.4mm,right=2.4mm,top=1.6mm,bottom=1.6mm
}

\begin{document}

% ==================== Page 1 ====================
\thispagestyle{empty}
\vspace*{0.08cm}
{\fontsize{22}{27}\selectfont\bfseries\color{navy}
Agent Skill Acquisition\\[-0.04cm]
文献与研究进展汇报\par}
\vspace{0.10cm}
{\large\color{slate}导师私聊汇报 · 组会选篇 · 研究定位修正}
\vspace{0.28cm}

\begin{tabularx}{\textwidth}{L{2.2cm}Y L{2.3cm}Y}
\toprule
\textbf{材料范围} & 已收集 44 篇 PDF；正文聚焦 31 篇高相关工作 &
\textbf{复核日期} & 2026-07-24 \\
\textbf{重点原文} & Semantic SC、SCR、HalluSquatting、SearchGEO、Skills That Don't Exist &
\textbf{已有实验} & Selection 三档 + acquisition runner 雏形 \\
\bottomrule
\end{tabularx}

\vspace{0.25cm}
\begin{keybox}
\textbf{汇报结论。} 不能再把贡献说成“首次研究 install”“首次测 E2E”或“首次测自主触发”。
当前更准确、可检验的问题是：\textbf{当用户只给普通能力任务、目标 skill 尚未安装时，agent
是否会自主完成外部发现、真实安装/注册和调用；这条链在模型输出、编排执行与系统状态之间具体断在哪里。}
\end{keybox}

\vspace{0.20cm}
{\large\bfseries\color{navy}阅读框架：把不同 endpoint 放回同一生命周期}

\begin{center}
\small
\fbox{\begin{minipage}{0.94\textwidth}
\centering
能力缺口 $\rightarrow$ 发起检索 $\rightarrow$ 检索命中 $\rightarrow$ 安装调用输出
$\rightarrow$ 解析/执行\\
$\rightarrow$ 文件落盘 $\rightarrow$ 注册可用 $\rightarrow$ 实际调用
$\rightarrow$ 载荷与任务结果
\end{minipage}}
\end{center}

\small
\begin{tabularx}{\textwidth}{L{2.65cm}L{4.45cm}Y}
\toprule
\textbf{文献区域} & \textbf{代表工作与可靠结论} & \textbf{对当前研究的影响} \\
\midrule
检索瓶颈 &
\textit{Overcoming the Retrieval Barrier}（USENIX Security 2026）：11 个数据集、8 个 embedding；Vanilla 主实验无法进入 top-5，CEM 后多设置接近满召回 &
提供“不要默认中间步骤自动发生”的方法模板 \\
\addlinespace
真实 skill 生态 &
\textit{Do Not Mention}（USENIX Security 2026）：98,380 $\rightarrow$ 4,287 $\rightarrow$ 157 个行为确认恶意 skill；632 个漏洞实例 &
证明恶意供给是真实生态问题，不只是合成 benchmark \\
\addlinespace
Registry 发现/选择 &
Semantic SC（\arxiv{2605.11418}）：Discovery 86.14\% pairwise、80\% Top-10；功能等价 paired selection 平均 77.6\% &
Discovery/selection 已有强邻居；需对齐竞争结构 \\
\addlinespace
装后危害 &
Skill-Inject / Poise：skill 已加载后可有高 payload ASR，并保持任务完成 &
post-load 不是主要 novelty，是 acquisition 后果上界 \\
\bottomrule
\end{tabularx}
\normalsize

\vspace{0.20cm}
\begin{warnbox}
\textbf{如何理解 funnel。} 它只是把“模型说要做”“runner 接受调用”“系统真实执行”“文件已经安装”
分开的测量框架。准确结论是：不同论文的\textbf{起点、终点和成功定义不同，数字不能直接横比}；
不能说“现有论文都把单段 ASR 冒充 E2E”。
\end{warnbox}

\newpage

% ==================== Page 2 ====================
\pageheading{1. 最近邻工作与研究边界}{原文复核后的诚实版本}

\scriptsize
\begin{tabularx}{\textwidth}{L{2.45cm}L{3.2cm}L{4.15cm}Y}
\toprule
\textbf{工作} & \textbf{触发/起点} & \textbf{实际测到} & \textbf{与待测问题的差异} \\
\midrule
Semantic SC\\\arxiv{2605.11418} &
skill 已在 registry 候选生命周期 &
Discovery、paired Selection、Governance 分段实验 &
不测真实安装、注册和后续调用的同轨迹 E2E \\
\addlinespace
SCR-CapFlow\\\arxiv{2606.15242} &
Neutral 中性任务语言；available skills &
observable mock state-change；Neutral 平均 33.6\%，单模型最高 91.5\% &
测 skill composition，不含外部发现/安装 \\
\addlinespace
SCR-TrustLift\\同上 &
上游 review/endorsement + 下游安装请求 &
simulated skill market 的有害安装；平均 1.10\% $\rightarrow$ 83.89\% &
不从普通任务自主发现开始；不测真实落盘/注册/调用 \\
\addlinespace
HalluSquatting\\\arxiv{2607.07433} &
用户明确要求 clone/install 资源 &
真实 agentic 应用中的 tool invocation / RCE；E2E 40\%--100\% &
触发已有 acquisition 意图，不是普通能力任务 \\
\addlinespace
SearchGEO\\\arxiv{2606.16821} &
信息/推荐任务；skill 部分为 18 例 probe &
GPT-5.4-mini 接受伪造 skill 并输出安装命令 17/18 &
停在 emitted command，不测执行/落盘 \\
\addlinespace
Skills That Don't Exist\\\arxiv{2607.12340} &
用户求 skill 推荐 &
主测 agent 幻觉率 36.9\%；明确指令下另有 benign install PoC &
不测 ordinary-task 自主恶意获取全链 \\
\bottomrule
\end{tabularx}
\normalsize

\vspace{0.22cm}
\begin{redbox}
\textbf{五条口径红线。}
\begin{itemize}
  \item 不说“第一个研究 install”：SCR、HalluSquatting、SearchGEO、Skills That Don't Exist 都已覆盖相邻端点。
  \item 不说“第一个测真实执行/E2E”：HalluSquatting 已测真实 RCE；SCR 用 observable state-change 计分。
  \item 不说“第一个测自主/中性触发”：SCR-CapFlow Neutral 已覆盖。
  \item 不说“SCR 的 skill 都是预装的”：这只适用于 CapFlow；TrustLift 已测 simulated market install。
  \item 不说“没人做过”：应说“在截至 2026-07-24 已核对的最接近工作中，尚未看到同时覆盖以下条件的研究”。
\end{itemize}
\end{redbox}

\vspace{0.18cm}
{\large\bfseries\color{navy}当前待检验的位置}

\begin{bluebox}
\textbf{Task-induced acquisition 的同轨迹系统测量：}
\begin{enumerate}
  \item 用户只给普通能力任务，不要求 search / install / skill；
  \item 目标 skill 在任务开始时尚未安装，agent 需从外部市场自主发现；
  \item 安装由真实 runner 解析并执行，用文件、manifest、registry 分别验证；
  \item 同一条轨迹继续测到调用、载荷与用户任务是否完成；
  \item 对每个阶段做可审计归因，并比较 model / scaffold / policy / tool budget。
\end{enumerate}
\end{bluebox}

\vspace{0.17cm}
\begin{warnbox}
\textbf{Research gap 本身不是贡献。} 只有当 pilot 得到跨任务/模型稳定的真实自主安装、
结构化 execution gap、强起点指标无法外推，或有效且低损的 acquisition gate，才可能形成顶会体量。
如果所有阶段都低且无结构，应按预先规则 No-Go。
\end{warnbox}

\newpage

% ==================== Page 3 ====================
\pageheading{2. 已有进展与下一步}{把先导观察和正式结论分开}

{\large\bfseries\color{navy}2.1 已完成：预装条件下的 Selection 实验}

\begin{tabularx}{\textwidth}{L{4.6cm}C{2.3cm}Y}
\toprule
\textbf{条件} & \textbf{结果} & \textbf{可支持的表述} \\
\midrule
Vanilla：朴素备份描述 & 0/20 & 本设定下未抢到调用 \\
Query+：通用“全能助手”描述 & 2/20 & 仅模糊/跨域任务偶发中招 \\
5 轮 LLM 对抗优化 & 1/20 & 当前优化方法未优于粗糙 baseline \\
\bottomrule
\end{tabularx}

\small
\textbf{实验条件：}HelloAgents + deepseek-chat；6 个良性专一 skill + 1 个恶意预装 skill；20 个自然任务。\\
\textbf{边界：}这是单模型、单框架、小样本的先导观察；不能写成“明确任务天然安全”，也不能与
Semantic SC 的功能等价 paired selection 77.6\% 直接比较。
\normalsize

\vspace{0.25cm}
{\large\bfseries\color{navy}2.2 已搭建：Acquisition 最小链路}

\begin{itemize}
  \item 已有 marketplace、\texttt{search\_skills}、\texttt{install\_skill} 和阶段日志雏形；
  \item 一次 A2 调试中，模型输出 \texttt{install\_skill(...)}，但默认 3 轮 tool budget 已耗尽，未进入真实执行；
  \item 这证明\textbf{指标必须拆分}，但还不能证明稳定的 execution gap。
\end{itemize}

\begin{center}
\small
\fbox{\begin{minipage}{0.94\textwidth}
\centering
\texttt{search\_called} $\rightarrow$ \texttt{target\_retrieved} $\rightarrow$
\texttt{install\_emitted} $\rightarrow$ \texttt{parsed} $\rightarrow$
\texttt{executed}\\
$\rightarrow$ \texttt{on\_disk} $\rightarrow$ \texttt{registered} $\rightarrow$
\texttt{invoked} $\rightarrow$ \texttt{payload} $\land$ \texttt{task\_ok}
\end{minipage}}
\end{center}

\vspace{0.20cm}
{\large\bfseries\color{navy}2.3 Pilot：先决定 Go / No-Go}

\begin{enumerate}
  \item \textbf{Instrumentation：}tool budget 可配置；完整 raw response / tool event；临时目录安装；
  \item \textbf{状态验证：}文件、manifest、registry、下一轮可调用分别验证；
  \item \textbf{任务条件：}ordinary / explicit-install / no-gap / irrelevant controls；
  \item \textbf{规模：}先做 60--100 runs，定位协议失败和阶段衰减；
  \item \textbf{冻结：}只有跨至少两个任务家族出现稳定结构，才扩大模型、scaffold 和 held-out tasks。
\end{enumerate}

\vspace{0.16cm}
\begin{keybox}
\textbf{可能值得继续的信号：}
\begin{itemize}
  \item ordinary task 下出现跨任务可复现的真实自主安装；
  \item emitted / executed / on-disk / registered 之间存在稳定落差，并能归因到 budget、parser、policy 或 scaffold；
  \item explicit-install 或 post-load 的高成功率不能预测 ordinary-task E2E，且差异可量化；
  \item 简单 acquisition gate / provenance policy 显著降风险且保持任务效用。
\end{itemize}
\end{keybox}

\newpage

% ==================== Page 4 ====================
\pageheading{3. 组会选篇与希望老师判断}{正式论文优先，最新邻居校准}

{\large\bfseries\color{navy}主推荐：两篇 USENIX Security 2026}

\begin{tabularx}{\textwidth}{L{4.8cm}Y}
\toprule
\textbf{论文} & \textbf{为什么值得精讲} \\
\midrule
\href{https://www.usenix.org/conference/usenixsecurity26/presentation/chang-hongyan}
{\textit{Overcoming the Retrieval Barrier}} &
问题定义、黑盒 CEM、规模实验、真实 E2E 与防御评估完整；是我从 selection barrier
走向 acquisition funnel 的方法学来源。重点学习“一个中间 barrier 如何被做成顶会论文”。 \\
\addlinespace
\href{https://www.usenix.org/conference/usenixsecurity26/presentation/liu-yi}
{\textit{“Do Not Mention This to the User”}} &
直接研究 agent skill 供应链；98,380 $\rightarrow$ 4,287 $\rightarrow$ 157 的测量漏斗、
动态沙箱、人工复核和 responsible disclosure 扎实。它回答“恶意 skill 是否真实存在”。 \\
\bottomrule
\end{tabularx}

\vspace{0.23cm}
\begin{bluebox}
\textbf{两篇的共同主线：}Do Not Mention 证明\textbf{恶意供给真实存在}；IPI 证明
\textbf{对象存在不等于能穿过中间获取门}。我的工作想测这两半在 agent skill acquisition 中如何接起来。
\end{bluebox}

\vspace{0.18cm}
{\large\bfseries\color{navy}最新邻居如何处理}

\begin{itemize}
  \item SCR、HalluSquatting 是定位上更近的必读工作，但截至 2026-07-24 公开状态为 arXiv 预印本；
  \item 组会用 1--2 页校准：SCR 已覆盖 Neutral composition 与 simulated install；
  HalluSquatting 已覆盖显式 install 条件下的真实 RCE；
  \item 若老师要求第二篇必须是“最贴当前 idea”，可用 \textbf{SCR 替换 Do Not Mention}。
\end{itemize}

\vspace{0.18cm}
{\large\bfseries\color{navy}希望老师重点判断}

\begin{enumerate}
  \item 是否认可把问题定位为\textbf{task-induced autonomous acquisition 的系统测量}，
  而不是“首次 install”或追一个孤立 ASR；
  \item 是否同意先用 60--100 runs pilot 判断阶段衰减有没有稳定结构，再决定是否扩大；
  \item 组会是否采用 \textbf{IPI + Do Not Mention}，并用 SCR / HalluSquatting 做最近邻比较；
  \item 若 pilot 通过，是否支持 L40 本地模型与后续一个外部模型做跨模型对照。
\end{enumerate}

\vfill
\small
\color{slate}
\textbf{安全说明：}本地实验仅使用合成任务与合成敏感数据；payload 只写 marker 或发送到
\texttt{127.0.0.1} collector，不进行真实外泄。完整文字版、组会提纲与私聊沟通稿见同目录
24 / 25 / 26 号 Markdown 文档。

\end{document}
